\documentclass[11pt,reqno]{amsart}
\usepackage[T1]{fontenc}
\usepackage[utf8]{inputenc}
\usepackage{lmodern}
\usepackage[a4paper,margin=27mm]{geometry}
\usepackage{amsmath,amssymb,mathtools,booktabs}
\usepackage{microtype}
\usepackage{xurl}
\usepackage[hidelinks]{hyperref}
\hypersetup{pdftitle={Polya inequalities for displaced square holes and perforated squares},pdfauthor={R. Arun Chandru}}
\numberwithin{equation}{section}
\setcounter{tocdepth}{1}
\allowdisplaybreaks[2]
\emergencystretch=2em
\raggedbottom
\newcommand{\R}{\mathbb R}
\newcommand{\N}{\mathbb N}
\newcommand{\Z}{\mathbb Z}
\newcommand{\dd}{\,\mathrm d}
\newcommand{\abs}[1]{\lvert#1\rvert}
\newcommand{\norm}[1]{\lVert#1\rVert}
\newcommand{\Om}{\Omega}
\newcommand{\sstar}{s_*}
\title[P\'olya inequalities for displaced square holes]{P\'olya inequalities for displaced square holes\texorpdfstring{\\}{ }and perforated squares}
\author{R. Arun Chandru}
\email{arun.chandru@gmail.com}
\date{28 September 2026}
\subjclass[2020]{Primary 35P15; Secondary 35P20, 11P21, 65G20}
\keywords{P\'olya's conjecture, Dirichlet Laplacian, displaced square hole, perforated domain, eigenvalue counting function, variational comparison, exact computation}
\begin{document}
\begin{abstract}
We prove P\'olya's inequality for every Dirichlet eigenvalue of a
square with an arbitrary strictly interior parallel square removed.
The hole may have any size and any position; neither centering nor
a positive lower bound on the boundary gaps is required. The proof
combines complementary square counts, endpoint-sensitive spectral
gaps, and a displacement-uniform transport between the actual
Dirichlet form domains. Mixed-boundary comparisons and explicit
scalar estimates cover the thin-gap regime. We also prove the
inequality for finitely many independently positioned and oriented
square holes whose total side length is at most
$\sqrt{1-12/(5\pi)}$ times the outer side. This threshold is exact
for the universal square-deficit subtraction criterion, not for
the geometric inequality. A further conclusion treats an arbitrary
parallel rectangular hole whenever both relative side lengths are
at least $11/12$, uniformly in their two gap scales and in position.
Analytic estimates control all unbounded spectral ranges. Exact
integer and rational certificates verify the remaining real
parameter intervals, without approximating eigenvalues of a
perforated domain.
\end{abstract}
\maketitle
\tableofcontents

\section{Introduction}\label{sec:introduction}

Weyl's law determines the leading density of the Dirichlet spectrum,
but does not decide on which side of its leading term the counting
function lies at a specified energy. P\'olya's conjecture asks for
this one-sided assertion at every energy. In the plane it states
that
\begin{equation}\label{eq:polya}
 |\Om|\lambda_k^D(\Om)\ge4\pi k\qquad(k\ge1),
\end{equation}
where the eigenvalues of the ordinary Euclidean Dirichlet Laplacian
are repeated according to multiplicity. The coefficient is forced
by Weyl's law. The difficulty is to retain it simultaneously through
the low spectrum, intermediate spectral clusters, and the unbounded
high-energy range.

Perforated domains expose a basic tension in this question.
Removing a hole increases Dirichlet eigenvalues, but decreases the
area multiplying them in \eqref{eq:polya}. Domain monotonicity
controls only the first change. Moreover, moving a hole destroys
symmetry without changing its area or its isolated spectrum. A
comparison that uses the isolated spectra alone is consequently
position-independent, but need not retain enough information to
prove the sharp inequality. This paper combines that comparison
with geometric form estimates and a transport of indexed spectral
gaps. The transport constants do not depend on the displacement.

\subsection*{The geometric conclusions}

Let $Q$ be an open square of side $L>0$, and let $H$ be a closed
square strictly inside $Q$, with sides parallel to those of $Q$.
Our principal conclusion is
\begin{equation}\label{eq:main-frame}
 \boxed{|Q\setminus H|\lambda_k^D(Q\setminus H)\ge4\pi k
       \qquad(k\ge1).}
\end{equation}
The side length and position of $H$ are arbitrary. In particular,
there is no centering assumption, no exceptional interval of hole
sizes, and no positive lower bound on the distance between boundary
components. Strict containment is required for each domain;
boundary contact is not included.

For an explicit parameterization, normalize $L=1$ and put
\begin{equation}\label{eq:frame}
\begin{split}
 Q&=(0,1)^2,\\
 H_{q,a,b}&=[(1-q)b,(1-q)b+q]\times
              [(1-q)a,(1-q)a+q],\\
 \Om_{q,a,b}&=Q\setminus H_{q,a,b},
 \qquad 0<q,a,b<1.
\end{split}
\end{equation}
Thus \eqref{eq:main-frame} is the assertion
$(1-q^2)\lambda_k^D(\Om_{q,a,b})\ge4\pi k$ for every value
of all three real parameters and every positive integer $k$.
Similarities preserve this area-normalized assertion.

Two further conclusions describe complementary geometries. For
finitely many pairwise disjoint closed squares $H_i$ strictly inside
$Q$, let $l_i$ denote their side lengths. Their positions and
orientations may be independent. Define
\begin{equation}\label{eq:sharp-constant}
 s_* =\sqrt{1-\frac{12}{5\pi}}
       =0.485856226839692\ldots.
\end{equation}
Then
\begin{equation}\label{eq:main-holes}
 \sum_i l_i\le s_*L
 \quad\Longrightarrow\quad
 \left|Q\setminus\bigcup_iH_i\right|
 \lambda_k^D\!\left(Q\setminus\bigcup_iH_i\right)\ge4\pi k
 \quad(k\ge1).
\end{equation}
The finite family of holes has strictly separated boundary
components. The sum-of-sides hypothesis, rather than a hypothesis
on the sum of their areas, is important here.

Finally, suppose that a single closed parallel rectangle $H$ has
side lengths $pL,qL$ and is strictly inside $Q$, at an arbitrary
position. We obtain
\begin{equation}\label{eq:main-displaced}
 \frac{11}{12}\le p,q<1
 \quad\Longrightarrow\quad
 |Q\setminus H|\lambda_k^D(Q\setminus H)\ge4\pi k
 \quad(k\ge1).
\end{equation}
The side ratios vary independently. The horizontal and vertical
gap totals need not be comparable, and each positive individual
gap can be arbitrarily small. The constant $11/12$ is sufficient,
not asserted optimal.

The precise sharpness statement in this paper concerns the scalar
comparison underlying \eqref{eq:main-holes}. With
\begin{equation}\label{eq:count-intro}
 C(t)=\#\{(m,n)\in\N^2:m^2+n^2\le t^2\},
 \qquad d(t)=\frac\pi4t^2-C(t),
\end{equation}
where $\N=\{1,2,\ldots\}$, we prove
\begin{equation}\label{eq:scalar-main}
 \sum_i d(q_it)\le d(t)
 \qquad\left(q_i\ge0,\ \sum_iq_i\le s_*,\ t\ge0\right).
\end{equation}
The number $s_*$ is the largest total-side threshold for which
this scalar assertion holds uniformly over finite lists.
This is a limitation of the comparison, not an optimal geometric
restriction. Indeed, \eqref{eq:main-frame} proves the geometric
inequality for single parallel square holes of every size.

\subsection*{Prior work and the contribution here}

P\'olya established \eqref{eq:polya} for tiling planar domains
\cite{Polya1961}. The inequalities of Berezin and of Li and Yau
control eigenvalue averages \cite{Berezin1972,LiYau1983}; they do
not directly imply the conjectured bound at each individual index.
Product constructions, packing arguments and subdivisions give
other classes and reductions
\cite{Laptev1997,FLP2021,FS2023,HW2026}. Filonov, Levitin,
Polterovich and Sher proved the conjecture for Euclidean balls
\cite{FLPS2023} and for circular planar annuli \cite{FLPS2026}.
The annulus argument combines variational bounds, special-function
estimates, lattice counting and rigorous finite computation.
Conformal equivalence does not transfer the ordinary Dirichlet
spectrum or its sharp area-normalized bound from circular annuli
to the polygonal domains considered here.

The closest comparison principle is the subtraction of
complementary counting functions used by Jiang and Lin for
perforated strip-tiling domains \cite{JL2026}. Their
Theorem~1.15, in dimension two, gives the sufficient condition
$\sum_i l_i/L\le1/(2\sqrt2\pi)$ for square holes in a square.
Their discussion also includes rectangular complementary counts
and the combination of low- and high-energy estimates. We use
this established variational principle. The contribution in
\eqref{eq:main-holes} is the exact threshold for the universal
square-deficit comparison, including the aggregation of an
arbitrary finite list of hole sizes.

The principal geometric contribution is the complete family
\eqref{eq:main-frame}, with both displacement parameters
unrestricted. Its proof supplements subtraction by an
endpoint-sensitive lower bound on indexed eigenvalues and a
piecewise-affine map between the actual form domains. The
resulting constants are independent of the two gap fractions.
For \eqref{eq:main-displaced}, a one-dimensional counting
comparison removes one displacement parameter, while a concavity
argument at low energy separates the other from two independent
gap scales.

All conclusions concern unsmoothed counts at every energy.
Weyl asymptotics, quantitative remainders and Riesz-mean
asymptotics provide important context
\cite{Weyl1912,JL2026,FrankLarson2025}, but do not alone resolve
the finite spectral range with the exact conjectured coefficient.
Our elementary disk discrepancy estimate supplies an explicit
tail. Rational interval and integer-lattice certificates settle
the remaining comparisons. No approximate eigenvalues of a
perforated domain enter the proof.

\subsection*{Organization of the argument}

Section~\ref{sec:comparison} develops the form comparisons and
the transport of spectral gaps, with strict and inclusive counting
conventions kept separate. Section~\ref{sec:sums} supplies the
one-dimensional estimates used when gaps are thin.
Section~\ref{sec:holes} proves the sharp total-side criterion.
Section~\ref{sec:all-squares} joins two continuous size bands to
the small- and large-hole ranges, proving
\eqref{eq:main-frame}. Section~\ref{sec:displaced} proves the
large rectangular-hole assertion independently of the disk
discrepancy estimate. The appendices give the analytic discrepancy
proof and every finite certificate obligation.

The finite computations are part of a proof over real parameter
intervals, not a grid of sampled domains. Each accepted interval
or box has a bound valid throughout it, its boundary is included
with the correct spectral convention, and an analytic estimate
controls the unbounded range. The scope is precisely
\eqref{eq:main-frame}, \eqref{eq:main-holes} and
\eqref{eq:main-displaced}. There is no assertion for arbitrary
planar domains, unrestricted large rotated holes, all rectangular
holes, or Neumann eigenvalues.

\section{Form comparisons and transported spectral gaps}\label{sec:comparison}

Every domain used below is bounded and polygonal, with finitely
many separated boundary components. Its Dirichlet Laplacian is
the selfadjoint operator associated with
\[
 \mathfrak a_\Om[u]=\int_\Om|\nabla u|^2\dd x,
 \qquad u\in H^1_0(\Om),
\]
in $L^2(\Om)$. The resolvent is compact, and
\begin{equation}\label{eq:minmax}
 \lambda_k^D(\Om)=
 \inf_{\substack{V\subset H^1_0(\Om)\\\dim V=k}}
 \sup_{0\ne u\in V}
 \frac{\int_\Om|\nabla u|^2}{\int_\Om|u|^2}.
\end{equation}
We use the inclusive counting function
\begin{equation}\label{eq:inclusive}
 N_\Om(E)=\#\{k:\lambda_k^D(\Om)\le E\}.
\end{equation}
Thus $N_\Om(E)\le|\Om|E/(4\pi)$ at every energy implies
\eqref{eq:polya} by evaluation at an eigenvalue. When a spectral
gap gives only a lower bound on an indexed eigenvalue, we use
that bound directly rather than infer an inclusive count at the
gap endpoint.

\subsection*{Inserting the hole boundaries}

Normalize the outer side to one and set
\begin{equation}\label{eq:energy-normalization}
 t=\frac{\sqrt E}{\pi},\qquad e=t^2.
\end{equation}
The eigenvalues of an isolated square of side $q$ are
$\pi^2(m^2+n^2)/q^2$, regardless of its position or orientation.
Let $P$ be the outer square with finitely many square holes
removed. Extension by zero embeds the direct sum of the
Dirichlet form domains on $P$ and the hole interiors into the
form domain on the outer square, preserving norms and energies.
Min--max gives
\begin{equation}\label{eq:subtraction}
 N_P(E)+\sum_i C(q_it)\le C(t).
\end{equation}
In particular, \eqref{eq:scalar-main} implies
\begin{equation}\label{eq:subtraction-weyl}
 N_P(E)\le C(t)-\sum_iC(q_it)
 \le\frac\pi4\left(1-\sum_iq_i^2\right)t^2
 =\frac{|P|}{4\pi}E.
\end{equation}
The inclusion of form domains, rather than an assumption on
eigenfunctions at the interfaces, fixes the direction of this
inequality.

\subsection*{Transport between actual displaced-hole operators}

Fix the normalized gap fractions $a,b\in(0,1)$ in
\eqref{eq:frame}. If $0<q\le r<1$, map each of the three
coordinate intervals (first gap, hole projection, last gap)
linearly from the $q$ configuration to the $r$ configuration.
The outer endpoints are fixed. The two gap slopes are the same,
independently of their relative lengths:
\[
 A=\frac{1-r}{1-q}\le1,\qquad C_0=\frac rq\ge1,
\]
where $C_0$ is the slope on the middle interval. The product map
$\Xi$ is a bi-Lipschitz bijection from $\Om_{q,a,b}$ onto
$\Om_{r,a,b}$. Its derivative is, almost everywhere, one of
\[
 \operatorname{diag}(A,A),\quad
 \operatorname{diag}(A,C_0),\quad
 \operatorname{diag}(C_0,A).
\]
The region of type $\operatorname{diag}(C_0,C_0)$ is precisely
the deleted square.

For $v\in H^1_0(\Om_{q,a,b})$, put $u=v\circ\Xi^{-1}$.
Change of variables gives
\[
 \int_{\Om_{r,a,b}}|u|^2\ge A^2\int_{\Om_{q,a,b}}|v|^2,
 \qquad
 \int_{\Om_{r,a,b}}|\nabla u|^2
 \le\frac{C_0}{A}\int_{\Om_{q,a,b}}|\nabla v|^2.
\]
Indeed, the transformed energy matrix is the identity on the
first region and has eigenvalues $C_0/A,A/C_0$ on each of
the other two. Composition is an isomorphism of the form
domains. Applying \eqref{eq:minmax} to corresponding
$k$-dimensional spaces proves
\begin{equation}\label{eq:transport}
 \lambda_k^D(\Om_{q,a,b})\ge T(q,r)\lambda_k^D(\Om_{r,a,b}),
 \qquad
 T(q,r)=
 \begin{cases}
 \displaystyle\frac{q(1-r)^3}{r(1-q)^3},&q\le r,\\[3pt]
 1,&q\ge r.
 \end{cases}
\end{equation}
For the second case, the hole grows by inclusion when its side
increases with $a,b$ fixed: its left endpoints decrease and its
right endpoints increase. Domain monotonicity proves that case.
In particular, for $r\downarrow q$, monotonicity and the first
case squeeze each indexed eigenvalue between
$\lambda_k^D(\Om_{q,a,b})$ and
$T(q,r)^{-1}\lambda_k^D(\Om_{q,a,b})$. This proves the right
continuity used next, including at multiple eigenvalues.

\subsection*{A strict outer count and an inclusive hole count}

For squared radii $s\ge0$ define
\[
 \mathcal C_{\le}(s)=C(\sqrt s),\qquad
 \mathcal C_{<}(s)=\#\{(m,n)\in\N^2:m^2+n^2<s\}.
\]
For $0<r<1$ and $\tau>0$, put
\begin{equation}\label{eq:gap-index}
 D(r,\tau)=\mathcal C_{<}(\tau)-\mathcal C_{\le}(r^2\tau).
\end{equation}
This integer is nonnegative, because $r^2\tau<\tau$.
The useful endpoint estimate is
\begin{equation}\label{eq:endpoint-gap}
 \lambda_{D(r,\tau)+1}^D(\Om_{r,a,b})\ge\pi^2\tau.
\end{equation}

To see why the inner count can include equality, choose
$r'\in(r,1)$ and enlarge the hole to side $r'$, keeping $a,b$
fixed. Every inner mode
counted by $\mathcal C_{\le}(r^2\tau)$ then has eigenvalue
strictly below $\pi^2\tau$. For all energies $E<\pi^2\tau$
sufficiently close to this endpoint, the subtraction inequality
gives
\[
 N_{\Om_{r',a,b}}(E)
 \le\mathcal C_{<}(\tau)-\mathcal C_{\le}(r^2\tau).
\]
Consequently its $(D(r,\tau)+1)$st eigenvalue is at least
$\pi^2\tau$. Let $r'\downarrow r$ and use the continuity
just proved. This establishes \eqref{eq:endpoint-gap} without
discarding a multiplicity at either boundary of the comparison.

Combining the gap and transport gives, for every reference
$r\in(0,1)$,
\begin{equation}\label{eq:indexed-gap}
 k\ge D(r,\tau)+1
 \quad\Longrightarrow\quad
 \lambda_k^D(\Om_{q,a,b})\ge\pi^2T(q,r)\tau.
\end{equation}
These are bounds on the actual operator. Neither the auxiliary
count nor a spectrum computed numerically is being substituted
for that operator.

\subsection*{Why interval endpoints suffice for transport}

For a fixed reference $r$, set $J_r(q)=(1-q^2)T(q,r)$.
On $q\le r$ it is a positive constant times
$q(1+q)/(1-q)^2$, whose derivative is
$(1+3q)/(1-q)^3>0$. On $q\ge r$ it is $1-q^2$, which
decreases. The expressions agree at $r$. Hence, on any closed
interval $[l,h]\subset(0,1)$,
\begin{equation}\label{eq:endpoint-minimum}
 \min_{l\le q\le h}J_r(q)=\min\{J_r(l),J_r(h)\}.
\end{equation}
This converts a rational endpoint test into a statement for
every real hole size in the interval. Uniformity in $a,b$
was already built into \eqref{eq:transport}.

\section{Square-root sums and mixed-boundary counts}\label{sec:sums}

The estimates for thin gaps use one-dimensional frequencies.
For $x\ge0$ and $h>0$, write
\begin{equation}\label{eq:rect-width-sums}
 A_h(x)=\sum_{\substack{j\ge0\\j+1/2\le hx}}
    \sqrt{x^2-\frac{(j+1/2)^2}{h^2}},\qquad
 R_h(x)=\sum_{\substack{j\ge1\\j\le hx}}
    \sqrt{x^2-\frac{j^2}{h^2}},\qquad R=R_1.
\end{equation}
Set $A_0=R_0=0$. At a frequency threshold the radical is
zero. The Neumann zero modes in a spectral count will be
retained separately, even at equality.

\subsection*{The integer-point sum}

Let $g_x(y)=\sqrt{x^2-y^2}$ on $[0,x]$, and put
$c_0=\sqrt{2/27}$. For $x\ge1$, set $j=\lfloor x\rfloor$
and $r=x-j$. Concavity gives the composite trapezoid estimate
\[
 R(x)\le\int_0^j g_x(y)\dd y-\frac x2+\frac{g_x(j)}2.
\]
The chord on $[j,x]$ lies below $g_x$, so its integral is at
least $r g_x(j)/2$. Since
$g_x(j)=\sqrt{r(2x-r)}\le\sqrt{2x}\sqrt r$, and the maximum
of $(1-r)\sqrt r$ on $[0,1]$ is $2/(3\sqrt3)$,
\begin{equation}\label{eq:R-upper}
 R(x)\le\frac\pi4x^2-\frac x2+c_0\sqrt x
 \qquad(x\ge1).
\end{equation}
At an integral $x$ the endpoint contribution is zero.
We will also use, for every $x\ge0$, the simpler
right-endpoint-sum inequality
\begin{equation}\label{eq:R-integral}
 R(x)\le\int_0^x g_x(y)\dd y=\frac\pi4x^2.
\end{equation}

\subsection*{The midpoint sum, uniformly down to small widths}

For $y\ge1/2$, let $w$ be the largest half-integer not
exceeding $y$ and put $r=y-w\in[0,1)$. Apply the trapezoid
inequality on $[1/2,w]$, use the decreasing initial
half-interval, and use the chord on $[w,y]$. This gives
\[
 A_1(y)-\frac\pi4y^2
 \le\frac{1-r}{2}\sqrt{y^2-w^2}
 \le c_0\sqrt y.
\]
The argument includes the case of one summand and a summand
at threshold. Below $1/2$ the sum vanishes, so this upper
bound remains valid. For $y\ge3$, the inequality
$c_0/\sqrt3\le4/25$ yields a bound linear in $y$.
The remaining closed interval $[1/2,3]$ is covered by the
123 rational intervals specified in
Appendix~\ref{app:displaced-certificate}. Thus
\begin{equation}\label{eq:rect-midpoint}
 A_1(y)\le\frac\pi4y^2+\frac4{25}y
 \qquad(y\ge0).
\end{equation}
Both bounds will be needed: the asymptotic one is stronger
at a wide gap, while the linear one is uniform as its width
shrinks. In particular, $A_h(x)=A_1(hx)/h$ gives
\begin{equation}\label{eq:midpoint-two-bounds}
 A_h(x)-\frac\pi4hx^2
 \le\min\left\{\frac4{25}x,c_0\sqrt{\frac{x}{h}}\right\}
 \qquad(h>0).
\end{equation}
One must not use the second estimate with an implicitly
bounded $h^{-1/2}$ when a displacement is unrestricted.

\subsection*{Removing one displacement parameter}

For a parallel square hole, put $\delta=1-q$ and
$x=\delta\sqrt E/\pi$. The two full-width horizontal
rectangles have heights $\delta a,\delta(1-a)$; the two
middle side rectangles have widths $\delta b,\delta(1-b)$
and length $q$. Relax matching across the horizontal cuts,
and relax the entire inward horizontal edge of each
horizontal rectangle to Neumann. Its other three edges
remain Dirichlet. On each side rectangle, the two width
ends remain Dirichlet and the two length ends are Neumann.
Restriction embeds the original form domain in the resulting
direct sum, preserving energy and norm. These relaxed
spectra therefore give an upper bound on the count.

\begin{figure}[ht]
\centering
\setlength{\unitlength}{0.72mm}
\begin{picture}(168,110)(-8,-5)
 \thicklines
 \put(0,0){\framebox(100,100){}}
 \put(20,20){\framebox(44,44){}}
 \thinlines
 \multiput(0,20)(4,0){5}{\line(1,0){2}}
 \multiput(64,20)(4,0){9}{\line(1,0){2}}
 \multiput(0,64)(4,0){5}{\line(1,0){2}}
 \multiput(64,64)(4,0){9}{\line(1,0){2}}
 \put(42,42){\makebox(0,0){$H$}}
 \put(50,10){\makebox(0,0){$\delta a$}}
 \put(50,82){\makebox(0,0){$\delta(1-a)$}}
 \put(10,42){\makebox(0,0){$\delta b$}}
 \put(82,42){\makebox(0,0){$\delta(1-b)$}}
 \put(106,43){\makebox(0,0)[l]{side length $q$}}
 \put(50,105){\makebox(0,0){outer side $1$}}
\end{picture}
\caption{The four-rectangle comparison for a displaced square
hole, with $\delta=1-q$. Dashed segments are artificial cuts.
Matching across them is relaxed, and the entire inward edge
of each full-width horizontal rectangle is also relaxed to
Neumann. The hole's two vertical sides remain Dirichlet.
The labels indicate widths or heights, not areas.}
\label{fig:rectangles}
\end{figure}

For the horizontal rectangles, sum the longitudinal
Dirichlet counts and bound each floor by its argument.
After multiplication by $\delta$, their contribution is at
most $A_a(x)+A_{1-a}(x)$. A transverse side frequency
$\nu$ contributes $1+\lfloor q\sqrt{x^2-\nu^2}/\delta\rfloor$,
where the one counts a longitudinal Neumann zero mode.
For every $s\ge0$,
\[
 \lfloor bs\rfloor+\lfloor(1-b)s\rfloor\le\lfloor s\rfloor.
\]
Furthermore, for $0\le\nu\le x$,
$\sqrt{x^2-\nu^2}=\int_\nu^x s/\sqrt{x^2-s^2}\dd s$.
Integrating the floor inequality against this nonnegative
weight proves $R_b(x)+R_{1-b}(x)\le R(x)$.
Keeping the zero-mode count gives
\begin{equation}\label{eq:square-mixed}
 \delta N_{\Om_{q,a,b}}(E)
 \le A_a(x)+A_{1-a}(x)+qR(x)+\delta\lfloor x\rfloor.
\end{equation}
This step removes $b$ without estimating a small gap by an
asymptotic expression with a nonuniform constant.
The same form construction, with two independent total gap
widths, will be used in Section~\ref{sec:displaced}.

\section{The exact total-side threshold for square holes}\label{sec:holes}

The argument for arbitrarily placed holes is entirely in terms
of the deficit $d$ in \eqref{eq:count-intro}. We first construct
an envelope that can aggregate any number of holes, and then
retain additional partition information where that envelope is
too coarse.

\subsection*{A deficit envelope with a monotone slope}

The exact reflection identity
\begin{equation}\label{eq:reflection}
 G(u)=1+4\lfloor u\rfloor+4C(u),\qquad
 G(u)=\#\{z\in\Z^2:\abs z\le u\},
\end{equation}
and the estimate proved in Appendix~\ref{app:circle} give
\begin{equation}\label{eq:deficit-tail}
 u-\frac34-\frac{11}{4}u^{2/3}\le d(u)
 \le u+\frac14+\frac{11}{4}u^{2/3}\qquad(u\ge2).
\end{equation}
In particular, for $u\ge1000$,
\begin{equation}\label{eq:thousand}
 \frac{2897}{4000}\le\frac{d(u)}u
 \le\frac{5101}{4000}<\frac43.
\end{equation}
The finite integer checks in Appendix~\ref{app:certificate}
complete the two bounds
\begin{equation}\label{eq:deficit-upper}
 d(u)\le\frac43u\quad(u\ge0),\qquad
 d(u)\le\frac{1309}{1000}u\quad(0\le u\le\sqrt{18}).
\end{equation}
These upper bounds include left limits at every lattice jump.

Define $F(0)=0$ and
\begin{equation}\label{eq:envelope}
 F(u)=
 \begin{cases}
 \pi u^2/4,&0<u<\sqrt2,\\
 \max\{\pi\sqrt2\,u/4,\ \pi u^2/4-1\},
        &\sqrt2\le u<\sqrt5,\\
 (1309/1000)u,&\sqrt5\le u<\sqrt{18},\\
 (4/3)u,&u\ge\sqrt{18}.
 \end{cases}
\end{equation}
The inequality $d(u)\le F(u)$ follows from the vanishing of $C$
below $\sqrt2$, the bound $C\ge1$ on the next range, and
\eqref{eq:deficit-upper}. Moreover, $F(u)/u$ is nonnegative
and nondecreasing. In the second range it is the maximum of a
constant and the increasing function $\pi u/4-1/u$.
At $\sqrt2$ there is no downward jump. At $\sqrt5$, the
left-limit prefix check in \eqref{eq:deficit-upper} gives
$(5\pi/4-1)/\sqrt5\le1309/1000$, and
$\pi\sqrt2/4<1309/1000$. At $\sqrt{18}$ the slope increases.
It follows also that $F$ is nondecreasing.

If $U=\sum_i u_i>0$, this monotone slope gives
\begin{equation}\label{eq:aggregate-envelope}
 \sum_i d(u_i)\le\sum_iF(u_i)
 \le\sum_i u_i\frac{F(U)}U=F(U).
\end{equation}
Zero summands and $U=0$ cause no difficulty. Consequently
$\sum_i d(q_it)\le F(\sstar t)$ under the hypothesis in
\eqref{eq:scalar-main}.

\subsection*{All but one energy band}

Put $R=243/500$. Rational enclosures of $\pi$ imply
\begin{equation}\label{eq:sstar-bracket}
 \frac{97}{200}<\sstar<R,
 \qquad \frac\pi4(1-\sstar^2)=\frac35.
\end{equation}
The squared transition energies
$2/\sstar^2,5/\sstar^2,18/\sstar^2$ belong respectively to
$(8,9),(21,22),(76,77)$. Write $C_n=C(\sqrt n)$.
For the first quadratic branch of \eqref{eq:envelope}, the
required inequality is $3t^2/5\ge C(t)$; for the second
quadratic branch it is $3t^2/5\ge C(t)-1$. The exact tests
\begin{equation}\label{eq:low-tests}
 3n\ge5C_n\quad(0\le n\le8),\qquad
 3n\ge5(C_n-1)\quad(8\le n\le21)
\end{equation}
establish them. The linear branch in the maximum is covered by
\begin{equation}\label{eq:low-linear-test}
 L_n:=333n-424C_n\ge0,
 \qquad(500L_n)^2\ge2n(243\cdot333)^2
 \quad(8\le n\le21).
\end{equation}
Indeed, division by $\pi/4$ reduces its inequality at $t=\sqrt n$
to $n-4C_n/\pi\ge\sstar\sqrt{2n}$. The left side is at
least $L_n/333$, and \eqref{eq:low-linear-test} proves the
stronger bound with $R$ in place of $\sstar$. All quantities
are checked nonnegative before squaring.

These are interval checks. On $n\le t^2<n+1$, the count is
$C_n$; the quadratic differences increase, and the linear
difference has derivative $\pi t/2-\pi\sqrt2\sstar/4>0$
for $t\ge\sqrt8$. The overlap at $n=8$ covers the first
noninteger transition.

For the two larger branches, the exact checks give
\begin{align}
 \frac{d(t)}t&\ge\frac{318087}{500000}
       =\frac{1309}{1000}R,
 &&\sqrt{21}\le t<\sqrt{77},\quad t^2\notin[53,54),
 \label{eq:middle-smallhole}\\
 \frac{d(t)}t&\ge\frac{81}{125}=\frac43R,
 &&t\ge\sqrt{76}.\label{eq:large-smallhole}
\end{align}
The finite tests proving these statements are listed in
Appendix~\ref{app:certificate}. Between successive integer
squared radii, $d(t)/t=\pi t/4-C_n/t$ increases. The infinite
part of \eqref{eq:large-smallhole} follows from
\eqref{eq:thousand}, since $2897/4000>81/125$.
The stated ranges overlap the transitions of $F$, so they
prove $d(t)\ge F(\sstar t)$ except possibly when
$53\le t^2<54$.

\subsection*{Retaining the partition at the exceptional band}

For every finite list $u_i\ge0$ satisfying
$\sum_i u_i\le B:=18/5$, we have
\begin{equation}\label{eq:partition}
 \sum_i d(u_i)\le\frac92.
\end{equation}
Here is the full argument. If every $u_i<\sqrt2$, then
\[
 \sum_i d(u_i)=\frac\pi4\sum_i u_i^2
 \le\frac{\pi\sqrt2}{4}B
 <\frac{1917}{452}<\frac92,
\]
using $\pi<355/113$ and $\sqrt2<3/2$. Otherwise retain one
actual member $v\ge\sqrt2$. Dropping the other counts and
summing their quadratic costs gives
\begin{equation}\label{eq:partition-function}
 \sum_i d(u_i)\le
 \frac\pi4\{v^2+(B-v)^2\}-C(v),\qquad\sqrt2\le v\le B.
\end{equation}
This does not suppose that only one member exceeds $\sqrt2$.
On the constant-count intervals ending at
$\sqrt5,\sqrt8,\sqrt{10},B$, the expression is convex.
The counts are respectively $1,3,4,6$, since $B^2<13$.
At an interior jump the value drops. It is therefore enough
to check the initial endpoint, the three subsequent left
limits, and $B$.

Using the positive lower radical bounds
$\sqrt2>7/5$, $\sqrt5>223/100$, $\sqrt8>14/5$,
$\sqrt{10}>79/25$, and $\pi<355/113$, the corresponding
upper bounds for \eqref{eq:partition-function} are
\begin{equation}\label{eq:partition-endpoints}
 \frac{2488}{565},\quad\frac{49973}{11300},\quad
 \frac{442}{113},\quad\frac{11349}{2825},\quad
 \frac{2361}{565}.
\end{equation}
All are below $9/2$. The radical bounds are verified by
squaring, and each endpoint comparison is rational. This
proves \eqref{eq:partition}.

For $53\le t^2<54$ the total scaled side satisfies
$\sum_i q_it\le\sstar t<R\sqrt{54}<18/5$.
Meanwhile $C(t)=37$, and
\begin{equation}\label{eq:53-close}
 d(t)\ge\frac{53\pi}{4}-37
 >\frac{53}{4}\frac{333}{106}-37
 =\frac{37}{8}>\frac92.
\end{equation}
Thus \eqref{eq:partition} closes the entire exceptional band.
Together with \eqref{eq:aggregate-envelope} and the preceding
comparisons, it proves \eqref{eq:scalar-main} for every real
$t\ge0$. Formula~\eqref{eq:subtraction-weyl} now proves
\eqref{eq:main-holes}.

\subsection*{The exact limitation of this scalar criterion}

If $S>\sstar$, choose
$\sstar<q<\min\{S,\sqrt{2/5}\}$. This is possible because
$\sstar<\sqrt{2/5}$. At $t=\sqrt5$, the inclusive count is
$C(t)=3$, while $C(qt)=0$. The inequality
$d(t)\ge d(qt)$ would require
$(5\pi/4)(1-q^2)\ge3$, or $q\le\sstar$, a contradiction.
Thus no larger threshold on the total side can guarantee
\eqref{eq:scalar-main}. The comparison itself, rather than the
geometry of the perforated domain, is what is sharp here.

\section{Every parallel square hole, at every position}\label{sec:all-squares}

We now prove \eqref{eq:main-frame}. The small-hole result covers
$q\le s_*$, and Section~\ref{sec:displaced} will cover
$q\ge11/12$. It remains to cover two closed, overlapping-endpoint
bands:
\[
 \frac{12}{25}\le q\le\frac45,
 \qquad
 \frac45\le q\le\frac{11}{12}.
\]
Since $12/25<97/200<s_*$, the first band overlaps the
small-hole range. We treat each band's unbounded energy range
analytically, leaving a finite collection of indexed inequalities
that the transported gaps certify uniformly in position.

\subsection*{A narrower disk-deficit corridor}

For every $u\ge201$, we have
\begin{equation}\label{eq:corridor}
 \frac89u\le d(u)\le\frac{10}{9}u.
\end{equation}
Appendix~\ref{app:certificate} gives the complete finite
integer tests. The explicit discrepancy estimate
\eqref{eq:deficit-tail} supplies the infinite tail: for
$u\ge15625=25^3$ it gives
\[
 \frac{d(u)}u\ge\frac{27811}{31250}>\frac89,
 \qquad
 \frac{d(u)}u\le\frac{17344}{15625}<\frac{10}{9}.
\]
The finite part counts all positive lattice pairs through
squared radius $15625^2$, in bounded-memory segments. It is
not a selection of radii at which the desired inequality happens
to hold. The tests include inclusive lower values and the left
limits needed for upper bounds at lattice jumps.

If $q\le4/5$ and $t\ge201/q$, both $t$ and $qt$ are in
\eqref{eq:corridor}, so
\[
 d(qt)\le\frac{10}{9}qt\le\frac89t\le d(t).
\]
Consequently square subtraction proves
\begin{equation}\label{eq:wide-tail}
 N_{\Om_{q,a,b}}(E)\le\frac{1-q^2}{4\pi}E
 \quad\left(q\le\frac45,\ E\ge\pi^2(201/q)^2\right),
\end{equation}
for every displacement. No estimate of an ellipse count is
required in this band.

\subsection*{A width-scale tail uniform in displacement}

For the second band, let $\delta=1-q$ and $x=\delta\sqrt E/\pi$.
Subtracting the desired Weyl term from
\eqref{eq:square-mixed} reduces the question to
\begin{equation}\label{eq:thin-defect}
 \begin{split}
 D_q(x,a)={}&A_a(x)+A_{1-a}(x)+qR(x)
                +(1-q)\lfloor x\rfloor\\
           &-\frac\pi4(1+q)x^2\le0.
 \end{split}
\end{equation}
At fixed $x,a$, its derivative in $q$ is
$R(x)-\lfloor x\rfloor-\pi x^2/4\le0$, by
\eqref{eq:R-integral}. It is enough to prove the inequality
at $q=4/5$, without discretizing the hole-size parameter.

For $z>0$, define the normalized sums
\[
 \Phi(z)=\sum_{\substack{j\ge0\\j+1/2\le z}}
           \sqrt{1-((j+1/2)/z)^2},\qquad
 \Psi(z)=\sum_{\substack{m\ge1\\m\le z}}
           \sqrt{1-(m/z)^2},
\]
and set $\Phi(0)=\Psi(0)=0$. Each is nondecreasing, including
at activation thresholds. Put $y=x\min(a,1-a)$, so
$0\le y\le x/2$. For $x>0$,
\begin{equation}\label{eq:normalized-defect}
 \frac{D_{4/5}(x,a)}x
 =\Phi(y)+\Phi(x-y)+\frac45\Psi(x)
  +\frac15\frac{\lfloor x\rfloor}{x}-\frac{9\pi}{20}x.
\end{equation}
Section~\ref{sec:sums} gives
\begin{align}
 \Phi(z)-\frac\pi4z&\le
        \min\left\{\frac4{25},\frac{c_0}{\sqrt z}\right\}
        &&(z>0),\label{eq:Phi-bound}\\
 \Psi(x)-\frac\pi4x&\le-\frac12+\frac{c_0}{\sqrt x}
        &&(x\ge1),\label{eq:Psi-bound}
\end{align}
where $c_0=\sqrt{2/27}$. At $z=0$, the corresponding
midpoint excess is zero.

First suppose $x\ge32$ and $8\le y\le x/2$. On this
interval $y^{-1/2}+(x-y)^{-1/2}$ decreases with $y$.
Using \eqref{eq:normalized-defect} and $\lfloor x\rfloor\le x$
therefore gives
\[
 \frac{D_{4/5}(x,a)}x
 \le c_0\left(\frac1{\sqrt8}+\frac1{\sqrt{24}}
                         +\frac4{5\sqrt{32}}\right)-\frac15<0.
\]
For an exact comparison, use $c_0<3/11$ and
$1/\sqrt8<71/200$, $1/\sqrt{24}<41/200$,
$1/\sqrt{32}<177/1000$. Their substitution gives
$2631/13750<1/5$.

Second, for every $0\le y\le x/2$ and $x\ge256$, use
the uniform $4/25$ bound for the narrower contribution
and $c_0/\sqrt{x-y}\le c_0\sqrt{2/x}$ for the wider one.
Then
\begin{equation}\label{eq:thin-analytic-tail}
 \frac{D_{4/5}(x,a)}x
 \le-\frac1{25}
       +\frac{c_0(\sqrt2+4/5)}{\sqrt x}<0.
\end{equation}
For example, $c_0<3/11$, $\sqrt2<10/7$ and
$\sqrt x\ge16$ make the right side less than
$-31/15400$. This estimate is uniform as either gap
fraction tends to zero.

Only the closed rectangle
\begin{equation}\label{eq:thin-compact}
 32\le x\le256,\qquad0\le y\le8
\end{equation}
remains. Certifying this rectangle, including points beyond
$y\le x/2$ if any, is sufficient. Here in fact $y\le8$
and $x\ge32$ already imply that restriction. Monotonicity
of $\Phi,\Psi$ gives exact upper bounds on full rational
boxes. Appendix~\ref{app:new-certificates} describes the
integer test and the complete subdivision. It accepts
865759 covering boxes, with none left unresolved. Thus
\eqref{eq:thin-defect} holds for $q\ge4/5$, $x\ge32$,
all real gap fractions, and hence
\begin{equation}\label{eq:thin-uniform-tail}
 N_{\Om_{q,a,b}}(E)\le\frac{1-q^2}{4\pi}E
 \quad\left(q\ge\frac45,\ E\ge\pi^2\frac{32^2}{(1-q)^2}\right).
\end{equation}

\subsection*{Reducing a high-energy count to finitely many indices}

Fix a closed size interval $[l,h]\subset(0,1)$ within one
of the two bands, and let $e_0$ be a squared-energy onset
valid for every size in it. We use
\begin{equation}\label{eq:energy-onsets}
 e_0=
 \begin{cases}
 (201/l)^2,&[l,h]\subset[12/25,4/5],\\
 32^2/(1-h)^2,&[l,h]\subset[4/5,11/12].
 \end{cases}
\end{equation}
Let $\pi_-=103993/33102<\pi<104348/33215=\pi_+$;
the rational enclosures are certified in
Appendix~\ref{app:certificate}. Put
\begin{equation}\label{eq:finite-index-cutoff}
 K(l,h)=\left\lfloor\frac{\pi_+}{4}(1-l^2)e_0\right\rfloor.
\end{equation}
Every index $k>K(l,h)$ already satisfies the desired
eigenvalue inequality. Indeed,
$4\pi k/(1-q^2)>\pi^2e_0$ for $q\in[l,h]$.
If $\lambda_k^D(\Om_{q,a,b})$ were smaller than that target,
an energy strictly between the target and
$\max\{\lambda_k^D(\Om_{q,a,b}),\pi^2e_0\}$ would have
count at least $k$, whereas the tail bound would make it
strictly less than $k$. This contradiction also deals
correctly with an eigenvalue equal to the onset energy.

\subsection*{Certifying the remaining indexed inequalities}

For integer $\tau\ge1$, \eqref{eq:gap-index} becomes
\[
 D(r,\tau)=C_{\tau-1}-C_{\lfloor r^2\tau\rfloor},
 \qquad C_n=C(\sqrt n).
\]
For each rational reference $r$, keep the largest available
integer $\tau\le250000$ whose gap index is at most $k$:
\begin{equation}\label{eq:gap-frontier}
 B_r(k)=\max\bigl(\{\tau\in\N:\tau\le250000,
                     D(r,\tau)+1\le k\}\cup\{0\}\bigr).
\end{equation}
This finite maximum need not be the optimal operator bound;
every candidate in it is justified by
\eqref{eq:indexed-gap}. With $S=2^{40}$ define
\begin{equation}\label{eq:integer-coefficient}
 \mu(l,h;r)=\left\lfloor
 \frac{S\pi_-}{4}
 \min\{(1-l^2)T(l,r),(1-h^2)T(h,r)\}\right\rfloor.
\end{equation}
The entirely integer test
\begin{equation}\label{eq:integer-acceptance}
 \mu(l,h;r)B_r(k)\ge Sk
\end{equation}
proves $(1-q^2)\lambda_k^D(\Om_{q,a,b})\ge4\pi k$
for every $q\in[l,h]$ and every $a,b\in(0,1)$.
This follows by multiplying \eqref{eq:indexed-gap} by
$(1-q^2)/(4\pi)$, using \eqref{eq:endpoint-minimum},
and rounding only in the direction that weakens the lower
bound. The certificate is thus on a full interval of
geometries, not at its endpoints alone.

The accompanying rational reference data cover the two bands
by adjoining closed intervals. Default references verify
most indices; the remaining ones have explicit additional
reference gaps. Every such reference, its count and its
integer comparison are checked anew by the verifier.
The full totals are
\begin{center}
\small
\begin{tabular}{@{}lrrrr@{}}
\toprule
Size band & Intervals & Index tests & Extra gaps & Largest $K$\\
\midrule
$[12/25,4/5]$ & 582 & 24288264 & 130 & 105989\\
$[4/5,11/12]$ & 739 & 9638177 & 81 & 18515\\
\bottomrule
\end{tabular}
\end{center}
Each index test refers to an entire size interval and all
displacements. There are no unverified intervals or indices.
The construction and independent verification obligations
are specified in Appendix~\ref{app:new-certificates}.

The finite indexed bounds and the two analytic tails now
prove \eqref{eq:main-frame} on $[12/25,11/12]$.
Section~\ref{sec:holes} supplies every $q\le s_*$, and
Section~\ref{sec:displaced} supplies every $q\ge11/12$.
The overlaps include all endpoints. This completes the
argument for every $0<q<1$, every strictly interior
position and every eigenvalue index.

\section{Large parallel rectangular holes with arbitrary displacement}\label{sec:displaced}

The last geometric conclusion allows the two sides of the hole to vary
independently. Let the outer square have side one and let the strictly
interior parallel rectangle have side lengths $p,q\in[11/12,1)$.
Its position is arbitrary. We prove
\[
 N_\Om(E)\le\frac{1-pq}{4\pi}E\qquad(E\ge0).
\]
Scaling and evaluation at each eigenvalue then give
\eqref{eq:main-displaced}. This argument does not use the
disk-discrepancy estimate.

Rotate the whole configuration by a quarter turn, if necessary, so
that $p\le q$, with $p$ horizontal. Put
\[
 \delta=1-q,\qquad c=\frac{1-p}{1-q}\ge1,
 \qquad x=\delta\frac{\sqrt E}{\pi},\qquad z=cx.
\]
The two horizontal gaps have lengths $\delta a$ and
$\delta(1-a)$, and the side gaps have lengths $c\delta b$ and
$c\delta(1-b)$, for independent $a,b\in(0,1)$. The width-indexed sums $A_h,R_h$ and $R=R_1$ are those
defined in \eqref{eq:rect-width-sums}.

Cut the frame into its two full-width horizontal rectangles and
two middle side rectangles. Relax matching at the cuts and the
entire inward horizontal edge of each horizontal rectangle to
Neumann. Its other three edges remain Dirichlet. The side
rectangles are Dirichlet across their widths and Neumann at both
ends of their length $q$. This enlarges the original Dirichlet
form domain, with norm and energy unchanged on restrictions.

The horizontal transverse frequencies are
$(j+1/2)/(\delta a)$ and $(j+1/2)/(\delta(1-a))$; summing
their longitudinal Dirichlet counts and bounding floors gives
$\delta^{-1}(A_a+A_{1-a})$. Each side transverse frequency
contributes one plus a floor, with the one counting the
longitudinal Neumann zero mode even at equality.
For every $s\ge0$,
\[
 \lfloor bs\rfloor+\lfloor(1-b)s\rfloor\le\lfloor s\rfloor.
\]
Integration against $s/\sqrt{x^2-s^2}$ gives
$R_b(x)+R_{1-b}(x)\le R(x)$, because a frequency $\nu\le x$
contributes $\sqrt{x^2-\nu^2}$. Applying this comparison to
total side width $c$, and retaining the zero modes, yields
\begin{equation}\label{eq:rect-mixed}
 \delta N_\Om(E)\le
 A_a(x)+A_{1-a}(x)+\frac q cR(cx)+\delta\lfloor cx\rfloor.
\end{equation}
Here $R_c(x)=R(cx)/c$. Since $1-pq=\delta(1+qc)$, the
desired Weyl expression, multiplied by $\delta$, is
$\pi(1+qc)x^2/4$.

\subsection*{An extremal reduction of the side lengths}

For fixed $c,x,a$, let
\begin{equation}\label{eq:rect-D}
 D=A_a(x)+A_{1-a}(x)+\frac{1-\delta}{c}R(cx)
       +\delta\lfloor cx\rfloor
       -\frac\pi4\bigl(1+(1-\delta)c\bigr)x^2.
\end{equation}
The decreasing right-endpoint sum satisfies $R(u)\le\pi u^2/4$.
Consequently
\[
 \frac{\partial D}{\partial\delta}
 =\lfloor cx\rfloor+\frac1c\left(\frac\pi4(cx)^2-R(cx)\right)
 \ge0.
\]
The condition $p\ge11/12$ gives $\delta\le1/(12c)$.
It therefore suffices to take
\begin{equation}\label{eq:rect-worst}
 \delta=\frac1{12c},\qquad p=\frac{11}{12},
 \qquad q=1-\frac1{12c}.
\end{equation}
This is a scalar comparison at fixed $x$; no geometric inclusion
is asserted.

For $x>0$ and $z=cx\ge x$, put
\begin{align}
 d_R(z)&=\frac\pi4z^2-R(z),\notag\\
 \mathcal H(x,a)&=\frac{A_a(x)+A_{1-a}(x)}x-\frac\pi4x,
       \label{eq:rect-HKL}\\
 \mathcal K(z,x)&=\left(1-\frac{x}{12z}\right)\frac{d_R(z)}z
                     -\frac{\lfloor z\rfloor}{12z},\notag\\
 \mathcal L(z)&=\frac{11}{12}\frac{d_R(z)}z
                     -\frac{\lfloor z\rfloor}{12z}.\notag
\end{align}
In \eqref{eq:rect-worst}, $D/x=\mathcal H-\mathcal K$.
Moreover $d_R\ge0$ and $x\le z$ give
\begin{equation}\label{eq:rect-KL}
 \mathcal K(z,x)\ge\mathcal L(z).
\end{equation}

\subsection*{Scalar estimates with explicit tails}

Section~\ref{sec:sums} proved the uniform midpoint estimate
\eqref{eq:rect-midpoint}, with its finite certificate in
Appendix~\ref{app:displaced-certificate}.
The scaling identity $A_h(x)=A_1(hx)/h$ now gives
\begin{equation}\label{eq:rect-H-global}
 \mathcal H(x,a)\le\frac8{25}\qquad(x>0,\ 0\le a\le1).
\end{equation}

The integer-point version of the same trapezoid argument is
\eqref{eq:R-upper}; its initial endpoint contributes $-y/2$.
Using that estimate and $\lfloor z\rfloor\le z$ gives
\begin{equation}\label{eq:rect-L-tail}
 \mathcal L(z)\ge\frac38-\frac{11}{12}\sqrt{\frac2{27}}z^{-1/2}
 \qquad(z\ge1).
\end{equation}
The finite certificates in Appendix~\ref{app:displaced-certificate}
and this tail prove
\begin{equation}\label{eq:rect-L-bounds}
 \mathcal L(z)\ge\frac7{25}\quad(z\ge1),\qquad
 \mathcal L(z)\ge\frac8{25}\quad(z\ge4).
\end{equation}
Their finite ranges are $[1,7)$ and $[4,21)$, respectively.
For the two tails, the exact positive differences are
\[
 7\left(\frac38-\frac7{25}\right)^2
 -\left(\frac{11}{12}\right)^2\frac2{27}
 =\frac{9061}{9720000},
\]
and
\[
 21\left(\frac38-\frac8{25}\right)^2
 -\left(\frac{11}{12}\right)^2\frac2{27}
 =\frac{12463}{9720000}.
\]
A further compact certificate gives
\begin{equation}\label{eq:rect-H-middle}
 \mathcal H(x,a)\le\frac4{15}
 \qquad(2\le x\le4,\ 0\le a\le1).
\end{equation}

\subsection*{Concavity removes displacement at low energies}

For $1\le x\le2$ there is a more precise structural bound:
\begin{equation}\label{eq:rect-low-structure}
 A_a(x)+A_{1-a}(x)
 \le\max\{A_1(x),\,2\sqrt{x^2-1}\}.
\end{equation}
If only one width contributes, monotonicity in width bounds its
sum by $A_1(x)$. If both contribute nonzero terms, a second
frequency in either width would require total scaled width
strictly greater than $1/2+3/2=2$. Hence each width contributes
only its first nonzero term. The function
$h\mapsto\sqrt{x^2-1/(4h^2)}$ is concave on its active
interval. The symmetric sum is therefore maximized at
$a=1/2$, giving the second term in
\eqref{eq:rect-low-structure}. Additional terms at $x=2$ have
value zero, so the endpoint is included.

Write
\[
 h_0(x)=2\sqrt{1-x^{-2}}-\frac\pi4x\qquad(1\le x\le2).
\]
Then \eqref{eq:rect-midpoint} and
\eqref{eq:rect-low-structure} imply
$\mathcal H\le\max\{4/25,h_0\}$. The last finite
certificate is the explicit comparison
\begin{equation}\label{eq:rect-center-comparison}
 h_0(x)\le\mathcal K(z,x)
 \qquad(1\le x\le2,\ x\le z<4).
\end{equation}
Only the two spectral scales remain; no displacement variable
occurs in this certificate.

\subsection*{Completion of the geometric proof}

If $z\ge4$, equations \eqref{eq:rect-H-global},
\eqref{eq:rect-L-bounds} and \eqref{eq:rect-KL} give
$\mathcal H\le\mathcal K$. Suppose $z<4$, so $x\le z<4$.
For $2\le x<4$, use $\mathcal H\le4/15<7/25\le
\mathcal L\le\mathcal K$. For $1\le x\le2$,
\eqref{eq:rect-center-comparison} treats the $h_0$ branch,
while $4/25<7/25\le\mathcal K$ treats the other branch.

If $0<x<1$, both horizontal widths cannot contribute nonzero
terms, so $A_a+A_{1-a}\le A_1(x)$ and
$\mathcal H\le4/25$. When $z\ge1$, the lower bound
$\mathcal K\ge7/25$ finishes the comparison. If $z<1$,
both $R(z)$ and $\lfloor z\rfloor$ vanish, giving
\[
 \mathcal K(z,x)=\frac\pi4z-\frac\pi{48}x
 \ge\frac{11\pi}{48}x.
\]
For $x<1/2$ one has $\mathcal H=-\pi x/4<0$; for
$1/2\le x<1$ one has $\mathcal K\ge11\pi/96>4/25$.
The case $x=0$ is the zero-energy count and requires no division.
The stated closed ranges cover $x=1/2,1,2,4$ and $z=1,4$;
the other integer $z$ boundaries are covered in the finite
certificate by the interval beginning there.

Thus $D\le0$ at every energy and for every allowed position
and pair of side lengths. Equation~\eqref{eq:rect-mixed}
proves the required inclusive count inequality and hence
\eqref{eq:main-displaced}. This rectangular-hole conclusion does not cover
arbitrary intermediate pairs of side lengths,
large rotated holes, or an arbitrary outer rectangle.

\appendix
\section{An explicit disk-discrepancy estimate}\label{app:circle}

We supply the analytic input to \eqref{eq:deficit-tail}, including
its constant and onset. The bound is deliberately elementary;
no improvement in the Gauss-circle discrepancy exponent is claimed.
With the inclusive count $G$ from \eqref{eq:reflection}, we need
\begin{equation}\label{eq:gauss}
 \abs{G(r)-\pi r^2}\le11r^{2/3}\qquad(r\ge2).
\end{equation}

\subsection*{The Bessel bound}

The classical integral representation and large-argument
asymptotics for $J_1$ will be used in their usual real
normalization \cite[\S\S10.9, 10.17]{DLMF}. They imply
\begin{equation}\label{eq:J1-bound}
 \abs{J_1(x)}\le\min\{x/2,x^{-1/2}\}\qquad(x>0),
\end{equation}
as follows. Poisson's representation is
\[
 J_1(x)=\frac{x}{\pi}\int_{-1}^{1}e^{ixt}\sqrt{1-t^2}\dd t,
\]
which gives the bound $x/2$. It also gives the second bound
for $0<x\le1$. For $1\le x\le2$, the alternating power
series has decreasing terms. Its lower two-term truncation
is nonnegative, and its upper three-term truncation gives
\[
 0\le J_1(x)\le\frac x2-\frac{x^3}{16}+\frac{x^5}{384}
 \le\frac x2-\frac{5x^3}{96}
 \le\frac4{3\sqrt5}<\frac35.
\]
Since $(3/5)\sqrt2<1$, this proves the remaining bound on
that interval.

For $x\ge2$, put $u(x)=\sqrt x\,J_1(x)$ and
$b(x)=1-3/(4x^2)$. The Bessel equation reads $u''+bu=0$.
Consequently
\[
 \mathcal E(x)=u'(x)^2+b(x)u(x)^2,
 \qquad\mathcal E'(x)=b'(x)u(x)^2\ge0.
\]
The standard function and derivative asymptotics give
$\lim_{x\to\infty}\mathcal E(x)=2/\pi$. Therefore
$b(x)u(x)^2\le2/\pi$. Since $b(x)\ge13/16$,
\[
 \abs{J_1(x)}^2\le\frac{32}{13\pi x}<\frac1x,
\]
completing \eqref{eq:J1-bound}.

\subsection*{Smoothing a disk without losing its boundary points}

Use the Fourier transform
$\widehat f(\xi)=\int_{\R^2}f(x)e^{-2\pi i x\cdot\xi}\dd x$.
For the indicator $\chi_R$ of a disk of radius $R$,
\[
 \widehat\chi_R(\xi)=\frac{R}{\abs\xi}J_1(2\pi R\abs\xi)
 \qquad(\xi\ne0).
\]
Let $\mu_\varepsilon=\chi_\varepsilon/(\pi\varepsilon^2)$,
where $0<\varepsilon<r$. The pointwise convolution comparisons
give
\begin{equation}\label{eq:disk-sandwich}
 \sum_{m\in\Z^2}(\chi_{r-\varepsilon}*\mu_\varepsilon)(m)
 \le G(r)\le
 \sum_{m\in\Z^2}(\chi_{r+\varepsilon}*\mu_\varepsilon)(m).
\end{equation}
The lower convolution is zero at an exterior tangent point;
the upper convolution equals one throughout the closed
radius-$r$ disk. Hence the inclusive convention in $G$
introduces no boundary error.

Write $c=1/(\pi\sqrt{2\pi})$. From \eqref{eq:J1-bound},
\[
 \abs{\widehat\chi_R(m)}\le\sqrt{\frac{R}{2\pi}}\abs m^{-3/2},
 \qquad
 \abs{\widehat\mu_\varepsilon(m)}
 \le\min\{1,c(\varepsilon\abs m)^{-3/2}\}.
\]
There are exactly $8k$ points of $\Z^2$ with supremum norm
$k$, and their Euclidean norm is at least $k$. Put
$K=c^{2/3}/\varepsilon=1/(2^{1/3}\pi\varepsilon)$.
The relevant summand is decreasing, so
\begin{align}
 &\sum_{m\ne0}\abs m^{-3/2}
             \min\{1,c(\varepsilon\abs m)^{-3/2}\}\notag\\
 &\quad\le8\sum_{k\ge1}k^{-1/2}\min\{1,(K/k)^{3/2}\}
 \le8\int_0^\infty y^{-1/2}\min\{1,(K/y)^{3/2}\}\dd y
 =24\sqrt K.\label{eq:lattice-shell}
\end{align}
The Fourier series of the periodized convolution is
absolutely convergent: the product decay is $O(\abs m^{-3})$.
The convolution is continuous and compactly supported.
Periodization and uniqueness of its absolutely convergent
Fourier series therefore justify Poisson summation here.
Its nonzero frequencies have total absolute value at most
\[
 \frac{24}{2^{2/3}\pi}\sqrt{R/\varepsilon}.
\]
Combining both sides of \eqref{eq:disk-sandwich}, and using
$R=r+\varepsilon$ for a common bound, yields
\begin{equation}\label{eq:disk-error}
 \abs{G(r)-\pi r^2}\le
 2\pi r\varepsilon+\pi\varepsilon^2
 +\frac{24}{2^{2/3}\pi}\sqrt{(r+\varepsilon)/\varepsilon}.
\end{equation}

Choose $\varepsilon=\tfrac12r^{-1/3}$. After division by
$r^{2/3}$, the right side is
\[
 \pi+\frac\pi4r^{-4/3}
 +\frac{24}{2^{2/3}\pi}\sqrt{2+r^{-4/3}}.
\]
For $r\ge2$, use $r^{-4/3}<2/5$, $\pi<22/7$,
$\pi>157/50$, $2^{2/3}>19/12$, and
$\sqrt{12/5}<31/20$. The expression is less than
\begin{equation}\label{eq:gauss-constant}
 \frac{121}{35}+\frac{22320}{2983}
 =\frac{1142143}{104405}<11.
\end{equation}
All radical comparisons follow by taking the appropriate
positive powers. This proves \eqref{eq:gauss} with the
explicit onset $r=2$.

Finally \eqref{eq:reflection} gives
\[
 d(u)=\lfloor u\rfloor+\frac14
             -\frac{G(u)-\pi u^2}{4}.
\]
The inequalities $u-1\le\lfloor u\rfloor\le u$ now give
\eqref{eq:deficit-tail}, including its different constants
$-3/4$ and $+1/4$.

\section{Exact arithmetic and the disk-deficit certificates}\label{app:certificate}

This appendix specifies the finite arithmetic inputs to
Sections~\ref{sec:holes} and~\ref{sec:all-squares}.
The analytic reductions in those sections and
Appendix~\ref{app:circle} are part of the proof; no finite
enumeration alone asserts an all-energy conclusion.

\subsection*{Rational constants and inclusive lattice counts}

The Machin identity
\[
 \pi=16\arctan(1/5)-4\arctan(1/239)
\]
and the alternating series
\[
 \arctan(1/a)=\sum_{j\ge0}\frac{(-1)^j}{(2j+1)a^{2j+1}}
\]
give rational enclosures. An even partial sum is a lower
bound, and adding the next positive term gives an upper
bound. Twelve terms for each arctangent suffice for both
\begin{equation}\label{eq:pi-bracket}
 \frac{333}{106}<\pi<\frac{355}{113},
 \qquad
 \pi_-:=\frac{103993}{33102}<\pi<\frac{104348}{33215}=: \pi_+.
\end{equation}
For completeness, tangent addition gives
$\tan(4\arctan(1/5)-\arctan(1/239))=1$ exactly.
The same series bounds place that angle between zero and
one, hence in $(0,\pi/2)$, fixing the branch $\pi/4$.

For an integer cutoff $M$, form
\begin{equation}\label{eq:count-array}
 r_n=\#\{(a,b)\in\N^2:a^2+b^2=n\},\qquad
 C_n=\sum_{j=0}^n r_j\quad(0\le n\le M).
\end{equation}
Enumerating unordered pairs and assigning weight two off
the diagonal and one on it retains all multiplicities.
Alternatively, enumerate every ordered pair directly.
The identity
\begin{equation}\label{eq:count-row}
 C_n=\sum_{a=1}^{\lfloor\sqrt n\rfloor}
                 \left\lfloor\sqrt{n-a^2}\right\rfloor
\end{equation}
provides direct row checks. All root floors here are exact
integer square roots. The small-hole verifier uses
$M=1000000$, for which $C_M=784387$.

\subsection*{Upper bounds use left limits; lower bounds use inclusive values}

Between consecutive integer squared radii,
$d(u)/u=\pi u/4-C_n/u$ increases. To bound it above, the
finite tests therefore use left limits at the next possible
jump. Put $U_n=355n-452C_{n-1}$. Whenever $U_n>0$, check
\begin{align}
 9U_n^2&\le1808^2n
      &&(1\le n\le1000000),\label{eq:cert-upper}\\
 (1000U_n)^2&\le(1309\cdot452)^2n
      &&(1\le n\le18).\label{eq:cert-prefix}
\end{align}
If $U_n\le0$ the upper comparison is immediate, without
squaring. These tests establish the finite parts of
\eqref{eq:deficit-upper}; \eqref{eq:thousand} covers every
$u\ge1000$, including the last endpoint.

For lower bounds put $L_n=333n-424C_n$, now using the
inclusive count. In addition to \eqref{eq:low-tests} and
\eqref{eq:low-linear-test}, check
\begin{align}
 L_n&\ge0,\quad
 (500000L_n)^2\ge(318087\cdot424)^2n
 &&(21\le n\le76,\ n\ne53),\label{eq:cert-hole-middle}\\
 L_n&\ge0,\quad
 (125L_n)^2\ge(81\cdot424)^2n
 &&(76\le n\le1000000).\label{eq:cert-hole-high}
\end{align}
The sole failure of the middle test with $n=53$ included
is exactly that index. The partition argument
\eqref{eq:partition}--\eqref{eq:53-close} supplies its whole
energy band. It is not omitted on numerical grounds.
Monotonicity between jumps, the overlapping transition
intervals in Section~\ref{sec:holes}, and
\eqref{eq:thousand} give the full real ranges.

The program \texttt{verify\_deficits.py} reconstructs these
counts and tests, the rational bounds placing $s_*$ and its
three transitions, the five partition endpoint bounds,
and the numerical constants in Appendix~\ref{app:circle}.
It takes no certificate data as input. Its arithmetic is
standard-library integer and rational arithmetic.

\subsection*{The corridor beginning at radius 201}

The program \texttt{verify\_corridor.py} constructs the
complete positive-quadrant count through
$M=15625^2=244140625$. Frequencies are accumulated in
segments of at most 4000000 unsigned 32-bit entries;
cumulative counts carry between adjacent segments.
Thus the frequency segment occupies at most 16 MB, not an
array proportional to the full cutoff.

On $n\le u^2<n+1$, the lower and upper endpoint tests
for \eqref{eq:corridor} are the conservative integer
inequalities
\begin{align}
 9(103993n-4\cdot33102C_n)
     &\ge32\cdot33102\lceil\sqrt n\rceil,
                  \label{eq:corridor-low-test}\\
 9(104348(n+1)-4\cdot33215C_n)
     &\le40\cdot33215\lfloor\sqrt{n+1}\rfloor.
                  \label{eq:corridor-high-test}
\end{align}
The first tests the inclusive lower endpoint, rounding the
required right side upward. The second tests the left limit
at the upper endpoint, rounding its permissible right side
downward. The use of rational lower and upper bounds for
$\pi$ has the same conservative directions.

The last lower failure occurs at $n=17226$, and the last
upper failure at $n=40320$. Both are below $201^2=40401$.
The final count is $C_{244140625}=191731890$.
All subsequent finite bands pass. The analytic tail in
Section~\ref{sec:all-squares} begins at radius 15625, so
there is no unverified endpoint between the finite and
infinite parts. The verifier explicitly requires the
certified onset to be at most 201.

For speed, this and the thin-tail verifier use NumPy
integer arrays. A floating square root is used only to
propose an integer $r$; the inequalities
$r^2\le n<(r+1)^2$ are then checked exactly before the
integer root is used. Failure of that validation stops
verification. Floating-point tolerances do not decide
any proof inequality. Array sizes, cumulative counts
and all fixed-width integer products are bounded before
use; scalar rational calculations use arbitrary-precision
Python integers.

\section{Interval certificates for displaced square holes}\label{app:new-certificates}

Two different finite tasks occur in the displaced-hole
proof: a scalar inequality in a compact two-dimensional
region, and indexed eigenvalue bounds on intervals of
hole sizes. Their analytic dependencies and their
acceptance tests are distinct.

\subsection*{The compact thin-tail rectangle}

Use denominator $D=2048$ and fixed-point scale $H=2^{20}$.
For $z=k/D>0$ and a positive integer $j$, put
\[
 V_{k,j}=4k^2-D^2j^2.
\]
When this is nonnegative, $\sqrt{V_{k,j}}/(2k)$ is a
summand of $\Phi(k/D)$ for odd $j$ and of $\Psi(k/D)$
for even $j$. The range $1\le j\le512$ contains every
possible summand needed for $0\le k/D\le256$.
An exact integer-root validation supplies
$U_{k,j}=\lceil\sqrt{V_{k,j}}\rceil$, with no added unit
when the radical is already a perfect square. Explicitly,
for $k>0$ put
\[
 f_k=\left\lceil\frac{H}{2k}
       \sum_{\substack{1\le j\le512\\j\text{ odd},\ V_{k,j}\ge0}}
       U_{k,j}\right\rceil,
 \qquad
 g_k=\left\lceil\frac{H}{2k}
       \sum_{\substack{1\le j\le512\\j\text{ even},\ V_{k,j}\ge0}}
       U_{k,j}\right\rceil.
\]
These integer arrays satisfy
\[
 \Phi(k/D)\le f_k/H,\qquad \Psi(k/D)\le g_k/H.
\]
Set $f_0=g_0=0$. The summation and division directions
are important; the arrays are upper enclosures, not
rounded approximations to the normalized sums.

Start with the 1792 closed unit boxes
$[n,n+1]\times[m,m+1]$, $32\le n\le255$, $0\le m\le7$.
For a box in integer coordinates
\[
 l/D\le x\le u/D,\qquad v/D\le y\le w/D,
\]
monotonicity in \eqref{eq:normalized-defect} gives
\begin{align*}
 5D_{4/5}(x,a)/x
 \le{}&5\Phi(w/D)+5\Phi((u-v)/D)+4\Psi(u/D)\\
      &+\frac{\lfloor u/D\rfloor}{l/D}
        -\frac{9\pi}{4}\frac lD.
\end{align*}
Consequently it suffices to check the integer inequality
\begin{equation}\label{eq:thin-box-test}
\begin{split}
 5f_w+5f_{u-v}+4g_u
 &+\left\lceil\frac{H\lfloor u/D\rfloor D}{l}\right\rceil\\
 &-\left\lfloor\frac{9\pi_-H l}{4D}\right\rfloor\le0.
\end{split}
\end{equation}
All positive terms have been rounded upward and the
subtracted term downward. The floor value at the right
endpoint includes a possible new mode, which is a safe
upper bound for the entire closed box.

If a box does not pass, split it at its two midpoints
into four closed children. Their union is their parent.
The subdivision must stop with a verification failure
if an unresolved box reaches the specified depth or
integer-grid limit; it is never silently discarded.
The completed run of \texttt{verify\_thin\_tail.py} has
\[
 \begin{array}{lr}
 \text{tested boxes}&1153748,\\
 \text{accepted boxes}&865759,\\
 \text{maximum subdivision level}&8,\\
 \text{unresolved boxes}&0.
 \end{array}
\]
The accepted integer areas sum to
$224\cdot8\cdot2048^2=7516192768$. Independently, the
tree identity is
$\text{tested}=1792+4(\text{tested}-\text{accepted})$.
Coverage follows from the actual subdivision procedure;
the area and tree identities are additional checks, not
substitutes for coverage of boundaries. The least
integer slack is zero, which is valid for the non-strict
inequality required here. The analytic comparisons for
the two infinite regions in
Section~\ref{sec:all-squares} are also checked as exact
rational inequalities by this verifier.

\subsection*{The rational reference data}

The two reference files are
\texttt{wide\_certificate.json} and
\texttt{thin\_certificate.json}. They specify the size
intervals and reference gaps for
\eqref{eq:integer-acceptance}. They are certificate
inputs, not accepted proof summaries. A reference size
is an integer numerator over 64000 in the first file
and over 192000 in the second. Both use squared-energy
cutoff 250000 and coefficient scale $S=2^{40}$.
The two denominators are not interchangeable.

For each interval $[l,h]$, the data record the cutoff
$K(l,h)$, two default reference numerators, their scaled
coefficients, the indices not covered by those two
references, and an additional reference and integer
energy for each such index. The verification does not
assume the accuracy of any of these entries. The
program \texttt{verify\_transport.py} does the following:
\begin{enumerate}
\item It certifies \eqref{eq:pi-bracket}, binds the
data to those rational constants and to the specified
scales, and checks that the ordered closed intervals
meet exactly, with the required first and last endpoints.
\item It reconstructs every ordered positive lattice
pair through squared radius 250000 and checks direct
row counts, including $C_{250000}=195837$.
\item It recomputes each cutoff using
\eqref{eq:energy-onsets}--\eqref{eq:finite-index-cutoff},
and every coefficient using
\eqref{eq:integer-coefficient} with exact fractions.
\item For each default reference $r$, it reconstructs
$B_r(k)$ from all $1\le\tau\le250000$ by assigning
each energy to its gap index and taking cumulative
maxima. Every index $1\le k\le K(l,h)$ is tested by
\eqref{eq:integer-acceptance}; the resulting uncovered
list must agree exactly with the data.
\item For each uncovered $k$, it verifies the additional
reference's gap index directly from
$C_{\tau-1}-C_{\lfloor r^2\tau\rfloor}+1$, verifies
that this index is at most $k$, and tests
$\mu(l,h;r)\tau\ge Sk$. The additional entries must
cover the uncovered list exactly once.
\end{enumerate}
Thus neither a saved success flag nor the program that
originally searched for references is a trusted premise.
The total is 1321 intervals, 33926441 interval/index
comparisons, and 211 additional gaps, with the bandwise
totals displayed in Section~\ref{sec:all-squares}.

All vectorized products are checked to fit signed
64-bit integers. For example,
$2^{40}\cdot250000<2^{63}$ and
$192000^2\cdot250000<2^{63}$ bound the two largest
types of scaled products. Coefficients are first
computed with arbitrary-precision fractions; the
verifier requires them to lie in $[0,2^{40})$ before
vector multiplication.

The transport verification alone supplies the finite
indexed part. The corridor in
Appendix~\ref{app:certificate} and the thin scalar
certificate above supply its respective high-energy
dependencies. Their combination with the form proofs,
the small-hole argument and the large rectangular-hole
argument is the proof of the complete geometric
assertion. It is not an extrapolation from finitely
many eigenvalue indices.

\section{Exact certificates for the rectangular-hole comparison}\label{app:displaced-certificate}

The accompanying program \texttt{verify\_rectangular.py}
verifies the five finite statements in Section~\ref{sec:displaced}.
It is independent of the other verifiers and uses only integer and
rational arithmetic. Twelve-term alternating arctangent enclosures
in the Machin identity recheck
$333/106<\pi<355/113$.

Put $P=333/106$ and $B_0=2^{24}$. For a nonnegative rational
number $r=u/v$, define
\begin{equation}\label{eq:rect-root-enclosure}
 U(r)=\operatorname{isqrt}\left(\left\lfloor
                  \frac{uB_0^2}{v}\right\rfloor\right)+1,
 \qquad T(r)=U(r)-1.
\end{equation}
Then $T(r)/B_0\le\sqrt r<U(r)/B_0$, also when $r=0$
or the radical is a perfect square. In the following acceptance
tests, every upper sum uses these enclosures after monotonicity
has moved its energy and width arguments to the indicated
endpoints. All failed comparisons cause subdivision; exhaustion
of a stated grid or depth is a fatal verification failure.

\subsection*{The midpoint bound and the lower envelopes}

For \eqref{eq:rect-midpoint}, use squared energy $r=y^2$ and
begin with $[1/4,9]$. A rational interval $[l,u]$ is accepted if
\begin{equation}\label{eq:rect-midpoint-box}
 \frac P4l+\frac4{25}\frac{T(l)}{B_0}
 \ge\operatorname{upper}A_1(\sqrt u).
\end{equation}
Otherwise bisect it. The 123 accepted intervals have total
width $35/4$, giving the full closed finite range
$1/2\le y\le3$.

For a lower bound $\mathcal L\ge\gamma$, divide the $z$
axis into half-open integer bands $n\le z<n+1$. On a
subinterval $[l,u)$ accept if
\begin{equation}\label{eq:rect-L-box}
 \frac{11P}{4}l^2
 \ge11\operatorname{upper}R(u)+n+12\gamma u.
\end{equation}
This is obtained by multiplying the desired inequality by
$12z$. Take $\gamma=7/25$ and $n=1,\ldots,6$, or
$\gamma=8/25$ and $n=4,\ldots,20$. Rational bisection
produces accepted widths summing to six and seventeen,
respectively. Their final endpoints are covered by
\eqref{eq:rect-L-tail}. At an integer boundary, the next
band uses the new floor value.

\subsection*{The remaining horizontal range}

For \eqref{eq:rect-H-middle}, let $r=x^2$ and use grid
denominators $G_0=K_0=2^{14}$ for $r$ and $a$. The initial
rectangle is $[4,16]\times[0,1/2]$; symmetry gives the
remaining displacement range. A box $[l,u]\times[a,b]$ is
accepted if
\begin{equation}\label{eq:rect-H-box}
 \frac P4l+\frac4{15}\frac{T(l)}{B_0}
 \ge\operatorname{upper}
       \{A_b(\sqrt u)+A_{1-a}(\sqrt u)\}.
\end{equation}
Both width sums increase in their width and squared energy,
so the separate endpoint maxima give a safe upper bound even
when they do not occur at a common displacement.

In integer coordinates $(r_l,r_u,a_l,a_u)$ the energy
coordinate is bisected when
$(r_u-r_l)K_0\ge2(a_u-a_l)r_u$, unless only the other
coordinate remains available. The accepted grid areas sum
exactly to $12G_0(K_0/2)=1610612736$.

\subsection*{The explicit two-scale comparison}

For $n=\lfloor z\rfloor$, subtraction gives
\begin{equation}\label{eq:rect-center-expanded}
 \mathcal K(z,x)-h_0(x)
 =\frac\pi4z+\frac{11\pi}{48}x
 -\left(1-\frac{x}{12z}\right)\frac{R(z)}z
 -\frac n{12z}-2\sqrt{1-x^{-2}}.
\end{equation}
The quotient $R(z)/z$ is a sum of terms
$\sqrt{1-j^2/z^2}$, each increasing in $z$ after activation.
The factor $1-x/(12z)$ increases in $z$ and decreases in $x$.
Therefore, on $[a,b]\times[c,d)$ within one fixed floor
band, it suffices to check
\begin{equation}\label{eq:rect-center-box}
 P\left(\frac c4+\frac{11a}{48}\right)
 \ge\left(1-\frac a{12d}\right)
       \operatorname{upper}\frac{R(d)}d
       +\frac n{12c}
       +2\frac{U(1-b^{-2})}{B_0}.
\end{equation}
Here the normalized Riesz sum is enclosed directly by summing
$U(1-j^2/d^2)/B_0$.

Use denominator $G_2=2^{16}$ for both coordinates and initial
rectangles $[1,2]\times[n,n+1)$, $n=1,2,3$. Intersect
each box with $z\ge x$. A box with zero intersection area is
discarded; all others are accepted by
\eqref{eq:rect-center-box} or subdivided along a longest
available coordinate, choosing $x$ at a tie.
The acceptance inequality proves the continuous frozen-$n$
expression on the full closed box and hence on its clipped
part within the half-open floor band. It does not substitute
the old floor value at the next integer.

For a grid rectangle $[x_l,x_u]\times[z_l,z_u]$, twice its
clipped grid area is computed exactly as
\[
 2(z_u-z_l)\max\{0,\min(x_u,z_l)-x_l\}
 +(v-u)(2z_u-v-u)
\]
when $v>u$, with the last term omitted otherwise, where
$u=\max(x_l,z_l)$ and $v=\min(x_u,z_u)$.
The sums of accepted doubled areas on the three bands are
$G_2^2,2G_2^2,2G_2^2$, exactly the corresponding initial
clipped areas.

Within each fixed floor band, diagonal points $z=x$ are
limits of the positive-area region above the diagonal.
The finite collection of accepted closed boxes therefore
covers them by continuity of the frozen-$n$ expression.
At $z=1,2,3$ the band beginning there supplies its actual
floor value; $z=4$ belongs to \eqref{eq:rect-L-bounds}.
Shared subdivision edges are included. Thus the area checks
are supplemented by a boundary argument and are not being
used alone to infer coverage of spectral jumps.

\subsection*{Replay totals}

The deterministic results are as follows.
\begin{center}
\small
\begin{tabular}{@{}lrrr@{}}
\toprule
Finite comparison & Tested & Accepted & Depth\\
\midrule
Midpoint, $1/2\le y\le3$ & 245 & 123 & 12\\
$\mathcal L\ge7/25$, $1\le z<7$ & 434 & 220 & 9\\
$\mathcal L\ge8/25$, $4\le z<21$ & 1995 & 1006 & 10\\
$\mathcal H\le4/15$, $2\le x\le4$ & 5395 & 2698 & 18\\
Center comparison, $1\le x\le2$, $x\le z<4$ & 9280 & 4620 & 18\\
\bottomrule
\end{tabular}
\end{center}
All 8667 accepted boxes have positive rational margins.
The triangle routine additionally skips 43 empty boxes.
The program checks these statistics and the interval-width
and grid-area identities against their exact expected values.
The least margins, in table order, are
\[
 \frac{501827}{4445962240},\quad
 \frac{39188871}{22229811200},\quad
 \frac{3232271}{4445962240},\quad
 \frac{39331}{6668943360},\quad
 \frac{440730641}{125696244449280}.
\]
These are enclosures on full intervals or boxes, not values
sampled at selected parameters. All unbounded ranges were
already reduced to the displayed finite regions by the
analytic estimates in Section~\ref{sec:displaced}.

Verification uses unconditional checks rather than
Python assertions that optimization flags can disable.
Caches are bounded. The computations certify explicit scalar
comparisons; no eigenvalue of a perforated domain is computed.


\section*{Note on AI Use}
The author declares the use of consumer-grade Gemini-3 and GPT-6 models for prior-art search, exploring and developing ideas and proofs, drafting and LaTeX typesetting; the author takes full and sole responsibility for all contents of this work.

\begin{thebibliography}{99}
\bibitem{Polya1961}
G. P\'olya,
On the eigenvalues of vibrating membranes,
\emph{Proc. London Math. Soc.} (3) \textbf{11} (1961), 419--433.
\href{https://doi.org/10.1112/plms/s3-11.1.419}{doi:10.1112/plms/s3-11.1.419}.

\bibitem{Berezin1972}
F. A. Berezin,
Covariant and contravariant symbols of operators,
\emph{Math. USSR-Izv.} \textbf{6} (1972), no.~5, 1117--1151.
\href{https://doi.org/10.1070/IM1972v006n05ABEH001913}{doi:10.1070/IM1972v006n05ABEH001913}.

\bibitem{LiYau1983}
P. Li and S.-T. Yau,
On the Schr\"odinger equation and the eigenvalue problem,
\emph{Comm. Math. Phys.} \textbf{88} (1983), no.~3, 309--318.
\href{https://doi.org/10.1007/BF01213210}{doi:10.1007/BF01213210}.

\bibitem{Laptev1997}
A. Laptev,
Dirichlet and Neumann eigenvalue problems on domains in Euclidean spaces,
\emph{J. Funct. Anal.} \textbf{151} (1997), no.~2, 531--545.
\href{https://doi.org/10.1006/jfan.1997.3155}{doi:10.1006/jfan.1997.3155}.

\bibitem{FLP2021}
P. Freitas, J. Lagac\'e, and J. Payette,
Optimal unions of scaled copies of domains and P\'olya's conjecture,
\emph{Ark. Mat.} \textbf{59} (2021), no.~1, 11--51.
\href{https://doi.org/10.4310/ARKIV.2021.v59.n1.a2}{doi:10.4310/ARKIV.2021.v59.n1.a2}.

\bibitem{FS2023}
P. Freitas and I. Salavessa,
Families of non-tiling domains satisfying P\'olya's conjecture,
\emph{J. Math. Phys.} \textbf{64} (2023), no.~12, 121503.
\href{https://doi.org/10.1063/5.0161050}{doi:10.1063/5.0161050}.

\bibitem{HW2026}
X. He and Z. Wang,
P\'olya's conjecture for thin products,
\emph{Int. Math. Res. Not. IMRN} (2026), no.~16, rnag182.
\href{https://doi.org/10.1093/imrn/rnag182}{doi:10.1093/imrn/rnag182};
\href{https://arxiv.org/abs/2402.12093v3}{arXiv:2402.12093v3}.

\bibitem{FLPS2023}
N. Filonov, M. Levitin, I. Polterovich, and D. A. Sher,
P\'olya's conjecture for Euclidean balls,
\emph{Invent. Math.} \textbf{234} (2023), no.~1, 129--169.
\href{https://doi.org/10.1007/s00222-023-01198-1}{doi:10.1007/s00222-023-01198-1}.

\bibitem{FLPS2026}
N. Filonov, M. Levitin, I. Polterovich, and D. A. Sher,
P\'olya's conjecture for Dirichlet eigenvalues of annuli,
\emph{J. London Math. Soc.} \textbf{113} (2026), no.~2, e70425.
\href{https://doi.org/10.1112/jlms.70425}{doi:10.1112/jlms.70425}.

\bibitem{JL2026}
R. Jiang and F. Lin,
P\'olya's conjecture up to $\epsilon$-loss and quantitative estimates
for the remainder of Weyl's law,
\emph{Comm. Pure Appl. Math.} (2026), e70058, early view.
\href{https://doi.org/10.1002/cpa.70058}{doi:10.1002/cpa.70058};
\href{https://arxiv.org/abs/2507.04307v4}{arXiv:2507.04307v4}.

\bibitem{Weyl1912}
H. Weyl,
Das asymptotische Verteilungsgesetz der Eigenwerte linearer partieller
Differentialgleichungen (mit einer Anwendung auf die Theorie der Hohlraumstrahlung),
\emph{Math. Ann.} \textbf{71} (1912), no.~4, 441--479.
\href{https://doi.org/10.1007/BF01456804}{doi:10.1007/BF01456804}.

\bibitem{FrankLarson2025}
R. L. Frank and S. Larson,
Riesz means asymptotics for Dirichlet and Neumann Laplacians on Lipschitz domains,
\emph{Invent. Math.} \textbf{241} (2025), 999--1079.
\href{https://doi.org/10.1007/s00222-025-01352-x}{doi:10.1007/s00222-025-01352-x}.

\bibitem{DLMF}
NIST Digital Library of Mathematical Functions,
release 1.2.8 of 2026-09-15,
\url{https://dlmf.nist.gov/}.
See \S\S10.9 and 10.17, in particular formulas 10.9.4, 10.17.3, and 10.17.9.
\end{thebibliography}
\end{document}
